\section{Curvatura, torsión y acción Einstein--Cartan--Holst}
\label{chap:gravedad-cartan-holst}

La estructura de transporte HMT distingue la variable de memoria, la
holonomía y el defecto de cierre. Su realización gravitatoria se formula
mediante una tétrada y una conexión. Esta sección establece la geometría de
primer orden y especifica el tipo del mapa entre ésta y el espacio de estados
TPK\@.

\subsection{Identidad discreta de Gauss–Stokes y transporte de incidencia a la frontera}

Sea \(K\) un complejo celular, sea \(W\) un grupo abeliano y sea
\(\Omega\in C^1(K;W)\) una cocadena. Denotamos por
\(\delta:C^1(K;W)\to C^2(K;W)\) el coborde celular y por
\(\langle\,\cdot\,,\,\cdot\,\rangle\) el emparejamiento entre cocadenas con
valores en \(W\) y cadenas celulares enteras.

\begin{theorem}[Identidad celular universal de Stokes]
Para toda \(2\)-cadena finita \(z\in C_2(K;\mathbb Z)\),
\[
 \boxed{\langle\delta\Omega,z\rangle
 =\langle\Omega,\partial z\rangle.}
\]
\end{theorem}
\begin{proof}
Es la identidad que define el coborde como operador adjunto de la frontera
celular. La finitud de \(z\) garantiza que ambos emparejamientos son sumas
finitas en \(W\).
\end{proof}

\begin{corollary}[Gauss--Stokes sobre una pseudovariedad celular]
Sea \(K\) una pseudovariedad celular regular, bidimensional, finita y compacta,
coherentemente orientada, posiblemente con borde. Si \(z_K\) es la suma de
sus \(2\)-celdas con la orientación coherente, entonces
\[
 \sum_{f\in K^{(2)}}(\delta\Omega)(f)
 =\sum_e[\partial z_K:e]\,\Omega(e).
\]
En particular, la regularidad permite identificar las incidencias no
canceladas de \(\partial z_K\) con las aristas del borde geométrico.
\end{corollary}
\begin{proof}
Se aplica el teorema a \(z_K\). La orientación coherente hace que cada arista
interior aparezca en \(\partial z_K\) dos veces con coeficientes opuestos;
las incidencias no canceladas forman la cadena de borde.
\end{proof}

La versión no abeliana no se obtiene sustituyendo sumas por productos sin más.
En esta obra se utiliza únicamente cuando se han fijado transportes a una base
común, un orden de las holonomías y la conjugación correspondiente. Fuera de
ese subdominio no se afirma una identidad de Stokes no abeliana.

\subsection{Capacidad areal de la celda qutrit: \(81\log 3\) y transporte colectivo de incidencia}

\begin{theorem}[Corte mínimo de un cubo discreto]
En el grafo cúbico \(N\times N\times N\) con capacidad unitaria, todo corte
que separa dos caras opuestas tiene capacidad al menos \(N^2\), y existe uno
de capacidad \(N^2\).
\end{theorem}
\begin{proof}
Hay \(N^2\) caminos disjuntos en aristas que atraviesan el cubo en la dirección
normal a las dos caras. Todo corte debe interceptarlos, luego su capacidad es
al menos \(N^2\). Una capa transversal contiene exactamente \(N^2\) aristas y
realiza la cota.
\end{proof}

Para \(N=9\), el corte vale 81 y
\[
 9^3=729=27^2.
\]
La última igualdad permite reagrupar seis trits como un volumen \(9^3\) o una
cuadrícula de cardinalidad \(27^2\). Es una biyección de conjuntos; no es un
isomorfismo entre
\((\ZZ/9\ZZ)^3\) y \((\ZZ/27\ZZ)^2\), cuyos exponentes son distintos.

\subsection{Tétrada, conexión y \(2\) defectos}

Sean \(e^I\) una tétrada y \(\omega^{IJ}=-\omega^{JI}\) una conexión de
Lorentz. Definimos
\[
 T_{\rm tor}^I=D_\omega e^I
 =\dd e^I+\omega^I{}_J\wedge e^J,
\]
\[
 F^{IJ}=\dd\omega^{IJ}+\omega^I{}_K\wedge\omega^{KJ},
 \qquad
 g=\eta_{IJ}e^I\otimes e^J.
\]
La descomposición
\[
 \omega^{IJ}=\mathring\omega^{IJ}(e)+K^{IJ}
\]
separa la conexión de Levi--Civita y la contorsión. Entonces
\[
 T_{\rm tor}^I=K^I{}_J\wedge e^J.
\]
Curvatura y torsión son defectos diferentes: \(F\) mide el cierre del
transporte de vectores; \(T_{\rm tor}\), el cierre de la tétrada. La memoria
\(T\) del principio de Ambivalencia tampoco es torsión geométrica: sólo un
funtor de realización puede relacionarlas.

\subsection{Acción de 1.er orden}

Sea \(M\) una variedad lisa orientada de dimensión cuatro. Si
\(\partial M\ne\varnothing\), suponemos variaciones de soporte compacto en el
interior o, equivalentemente para el problema variacional considerado, un
término de borde y condiciones de contorno que anulen la contribución
fronteriza. Cuando \(\Psi\) contenga campos espinoriales, suponemos además una
estructura de espín y campos compatibles con ella. Sean
\[
 \Sigma^{IJ}=e^I\wedge e^J,
 \qquad
 (\star X)^{IJ}=\frac12\epsilon^{IJ}{}_{KL}X^{KL},
 \qquad
 \kappa=\frac{8\pi G}{c^4},
 \qquad
 \gamma\in\mathbb R\setminus\{0\}.
\]
La acción Palatini--Cartan--Holst es
\[
\boxed{
 S[e,\omega,\Psi]
 =\frac1{2\kappa c}\int_M
 \Sigma_{IJ}\wedge\left(\star F^{IJ}+\frac1\gamma F^{IJ}\right)
 -\frac\Lambda{\kappa c}\int_M\operatorname{vol}_e
 +S_{\rm m}[e,\omega,\Psi].}
\]
La variación de esta acción determina las correcciones torsionales
\cite{Sciama1964,Kibble1961,Hehl1976,Holst1996}.

Definimos la corriente de espín por
\[
 \delta_\omega S_{\rm m}
 =-\frac12\int_M\delta\omega^{IJ}\wedge\sigma_{IJ}.
\]

\begin{theorem}[Ecuación de Cartan--Holst]
Bajo las hipótesis geométricas y de borde anteriores, la variación respecto
de \(\omega\) da
\[
 D_\omega\!\left[(\star+\gamma^{-1}I)\Sigma^{IJ}\right]
 =\kappa c\,\sigma^{IJ},
\]
con la normalización global fijada por la definición anterior de
\(\sigma^{IJ}\).
\end{theorem}
\begin{proof}
Se usa \(\delta F^{IJ}=D_\omega\delta\omega^{IJ}\), se integra por partes y
se iguala a cero el coeficiente de la variación arbitraria
\(\delta\omega^{IJ}\). Con la convención de signo fijada en la definición de
\(\sigma^{IJ}\), el término de materia pasa al miembro derecho con signo
positivo. El factor \(1/(2\kappa c)\) produce el coeficiente
\(\kappa c\). El término cosmológico no depende de \(\omega\), y las
condiciones de borde eliminan el término generado por la integración por
partes.
\end{proof}

En firma lorentziana \(\star^2=-I\) sobre bivectores y, para
\(\gamma\in\RR\setminus\{0\}\),
\[
 (\star+\gamma^{-1}I)^{-1}
 =\frac{\gamma^2}{1+\gamma^2}(-\star+\gamma^{-1}I).
\]

\begin{corollary}[Desacoplamiento torsional]
Para una co-tétrada no degenerada y
\(\gamma\in\mathbb R\setminus\{0\}\),
\[
 \sigma^{IJ}=0\Longrightarrow T_{\rm tor}^I=0
 \Longrightarrow\omega=\mathring\omega(e).
\]
\end{corollary}

Para hacer explícita la normalización, definimos las fuentes normalizadas por
la acción mediante
\[
 \delta_g S_{\rm m}
 =-\frac12\int_M\operatorname{vol}_e\,
 \mathcal T^{S}_{\mu\nu}\,\delta g^{\mu\nu}.
\]
Cuando las coordenadas de la tétrada tienen dimensión de longitud, el tensor
energía--momento físico convencional es
\(T^{\rm phys}_{\mu\nu}:=c\,\mathcal T^{S}_{\mu\nu}\). Denotamos de igual
modo por \(\mathcal U^{S,\rm spin}\) la contribución normalizada por la
acción y por \(U^{\rm phys,\rm spin}:=c\,\mathcal U^{S,\rm spin}\) su
convención física. Al eliminar algebraicamente la conexión, la ecuación
métrica adopta entonces las dos formas equivalentes
\[
 G_{\mu\nu}+\Lambda g_{\mu\nu}
 =\kappa c\,
 \bigl(\mathcal T_{\mu\nu}^{S,\rm LC}
       +\mathcal U_{\mu\nu}^{S,\rm spin}\bigr)
 =\kappa\,
 \bigl(T_{\mu\nu}^{\rm phys,LC}
       +U_{\mu\nu}^{\rm phys,spin}\bigr),
\]
donde la contribución de espín es cuadrática en la corriente correspondiente
y depende del acoplamiento de materia. El coeficiente de esa contribución no
se congela sin especificar la acción fermiónica y las convenciones de
dualidad.

\subsection{Identidad de Bianchi con torsión y balance de flujo en la conexión de Cartan}

Para toda conexión,
\[
 D_\omega T_{\rm tor}^I=F^I{}_J\wedge e^J,
 \qquad
 D_\omega F^{IJ}=0.
\]
Después de eliminar torsión,
\[
\nabla^\mu(T_{\mu\nu}^{\rm phys,LC}
          +U_{\mu\nu}^{\rm phys,spin})=0.
\]
Estas identidades son pruebas de consistencia obligatorias: cualquier término
gravitatorio HMT debe surgir de una acción compatible o satisfacer el balance
de Noether correspondiente.

\subsection{Mapa de realización y especialización del parámetro de Holst}

La inserción definida en el \cref{mdg:chap:barbero} toma
\[
 \gamma=\Gamma_{6\to8}.
\]
Una vez realizada esta sustitución, todas las consecuencias variacionales
anteriores son teoremas de la acción especializada. La selección física de
esta acción por el transporte HMT se expresa mediante un mapa de realización
\[
 \mathfrak R_{\rm grav}:\mathcal X_{\rm TPK}
 \dashrightarrow\{(e,\omega,\Psi)\}
\]
dotado de tres propiedades: entrelazar el transporte TPK con el de
\(\omega\), identificar las clases de holonomía y suministrar las escalas de
\(G\) y \(\Lambda\). Estas propiedades fijan el dominio preciso de una
realización gravitatoria del formalismo.

Para co-tétrada no degenerada, \(\gamma\) real no nulo y materia sin corriente
de espín, el sector torsional se reduce a la formulación de primer orden de
relatividad general con la misma constante cosmológica, sujeto a las
condiciones de borde ya declaradas. En un fluido de espín, el escalado
\(s^2\propto a^{-6}\) puede producir un rebote Einstein--Cartan; obtener su
densidad y sus condiciones iniciales a partir del estado TPK es una cuestión
distinta de la consistencia variacional de la acción.
